\clearpage
\section{Frozen fast energy defect}\label{b-frozen-fast-energy-defect}

Technical derivation for review. Equation and subsection numbers in this
appendix are local to the source derivation. Historical local labels are
retained, including numbering gaps or nonsequential labels. The
accompanying source notes retain the original expressions for
comparison.

This appendix writes the frozen energy in physical units. Multiply
\(E_{0}\) by \(2/R_{\mathrm{BFSS}}\) to use the kinetic-one form of the
main text; set \(\ell =1\) after the stated fixed dilation.

\subsection{1. Statement and
normalization}\label{b-statement-and-normalization}

Let n,m be positive integers, \(d=nm\), and let \(A_{i}\) and
\(B_{i}, i=1,\ldots ,9\), be Hermitian internal matrices of the two
blocks. Rotate the spatial coordinates so that their center difference
is \(b=r e_{9}, r>0\). On \(\operatorname{Hom} (C^{m},C^{n})\), with its
Hilbert--Schmidt metric, put

% DISPLAY b.1
% SOURCE T_i Z=A_i Z−Z B_i,
% SOURCE Q_i=r δ_i9 I_d+T_i,
% SOURCE H=sum_i Q_i²,     C_ij=[Q_i,Q_j]=[T_i,T_j].
\begin{gather*}
\begin{aligned}
& T_{i} Z=A_{i} Z-Z B_{i}, \\
& Q_{i}=r \delta _{i9} I_{d}+T_{i}, \\
& H=\sum _{i} Q_{i}^{2}, \quad  C_{ij}=[Q_{i},Q_{j}]=[T_{i},T_{j}].
\end{aligned}\notag
\end{gather*}

Assume each \(\Vert T_{i}\Vert _{\mathrm{op}}\le r/100\). A
coordinate-invariant sufficient condition is

% DISPLAY b.2
% SOURCE (sum_i ||A_i||F²)^(1/2)+(sum_i ||B_i||F²)^(1/2)≤r/100.
\begin{gather*}
\begin{aligned}
& (\sum _{i} \Vert A_{i}\Vert _{F}^{2})^{1/2}+(\sum _{i} \Vert B_{i}\Vert _{F}^{2})^{1/2}\le r/100.
\end{aligned}\notag
\end{gather*}

Thus it suffices that each endpoint radius is at most r/200. No relation
to any unrelated pair\textquotesingle s separation is imposed. Each pair
may use its own orthonormal spatial frame; a common global axial
direction for all pairs is unnecessary. All traces below are finite
ordinary, unnormalized traces.

On the slice \(Z_{9}=0\) define 8-by-8 matrices with d-by-d blocks:

% DISPLAY b.3
% SOURCE G_ij=δ_ij I_d+Q_i Q_9^(-2) Q_j,
% SOURCE K_ij=δ_ij H−Q_i Q_j+2 C_ij,       i,j≤8.
\begin{gather*}
\begin{aligned}
& G_{ij}= \delta _{ij} I_{d}+Q_{i} Q_{9}^{-2} Q_{j}, \\
& K_{ij}= \delta _{ij} H-Q_{i} Q_{j}+2 C_{ij}, \quad  i,j\le 8.
\end{aligned}\notag
\end{gather*}

The kinetic-aware bosonic frequency matrix is \(B=G^{1/2} K G^{1/2}\).
For Hermitian Spin(9) matrices satisfying
\(\{ \Gamma _{i}, \Gamma _{j}\}=2 \delta _{ij} I_{16}\) put

% DISPLAY b.4
% SOURCE D=sum_i Γ_i⊗Q_i.
\begin{gather*}
\begin{aligned}
& D=\sum _{i} \Gamma _{i} \otimes  Q_{i}.
\end{aligned}\notag
\end{gather*}

Both \(K\) and \(B\) are positive, and \(D\) has precisely 8d negative
and 8d positive eigenvalues. The frozen ground energy, in physical
units, is

% DISPLAY b.5
% SOURCE E_0=(R/ell³)[Tr_(8d) sqrt(B)−(1/2)Tr_(16d)|D|].
\begin{gather*}
\begin{aligned}
& E_{0}=(R/\ell ^{3})[\operatorname{Tr} _{8d} \sqrt{B}-(1/2)\operatorname{Tr} _{16d}\vert D\vert ].
\end{aligned}\notag
\end{gather*}

The following explicit, deliberately non-sharp constant is valid for
every n,m:

% DISPLAY b.6
% SOURCE |E_0| ≤ 150002 (R/ell³) r^(-3)
% SOURCE               sum_(i<j)||[T_i,T_j]||F².                    (1)
\begin{gather*}
\begin{aligned}
& \vert E_{0}\vert  \le  150002 (R/\ell ^{3}) r^{-3} \\
& \sum _{i<j}\Vert [T_{i},T_{j}]\Vert _{F}^{2}.
\end{aligned}\tag{1}
\end{gather*}

The constant in (1) is independent of \(d\); endpoint weights below
carry the finite-rank dependence. The proof neither diagonalizes
internal matrices nor invokes an internal Hamiltonian gap, internal
eigenvalue gap, regularity of the commuting variety, or a
nearest-commuting-tuple estimate.

\subsection{2. Why these are the correct frozen
matrices}\label{b-why-these-are-the-correct-frozen-matrices}

The eliminated off-block Gauss principal equation is

% DISPLAY b.7
% SOURCE Q_9 p_9+sum_(i≤8)Q_i p_i=0.
\begin{gather*}
\begin{aligned}
& Q_{9} p_{9}+\sum _{i\le 8}Q_{i} p_{i}=0.
\end{aligned}\notag
\end{gather*}

Its positive kinetic square gives \(G\) exactly. It must not be dropped
or treated as a separate small bounded perturbation.

For
\(V(X)=(R/(2\ell ^{6}))\sum _{i<j}\Vert [X_{i},X_{j}]\Vert _{F}^{2}\),
the quadratic part in one off-block pair is \((R/\ell ^{6}) Z^{*} K Z\).
The square of the linear off-block commutator gives
\(\delta _{ij} H-Q_{j} Q_{i}\). The internal-background commutator
paired with \([Y_{i},Y_{j}]\) contributes \(C_{ij}\). Their sum is

% DISPLAY b.8
% SOURCE δ_ij H−Q_j Q_i+C_ij=δ_ij H−Q_i Q_j+2C_ij.
\begin{gather*}
\begin{aligned}
& \delta _{ij} H-Q_{j} Q_{i}+C_{ij}= \delta _{ij} H-Q_{i} Q_{j}+2C_{ij}.
\end{aligned}\notag
\end{gather*}

This works for two arbitrary internal blocks. The second block
contributes with the minus sign in \(T_{i}\), including in its induced
commutator. At quadratic order other block pairs do not couple into this
pair\textquotesingle s Hessian. Interactions involving three distinct
blocks begin in the higher-degree cross terms and are not estimated
here.

The 8d complex bosonic coordinates are 16d real oscillators. Their
ground energy is \((R/\ell ^{3})\operatorname{Tr}  \sqrt{B}\). The fast
fermionic sector has 16d complex annihilators; filling the 8d negative
eigenmodes gives \(-(R/(2\ell ^{3}))\operatorname{Tr}  \vert D\vert\).
The reverse block is the adjoint, not an independent second copy.

\subsection{3. Uniform cross gaps}\label{b-uniform-cross-gaps}

Scale \(r\) to one and write \(\kappa  =1/100\). Then
\(Q_{9}=I+T_{9}, Q_{i}=T_{i}\) for \(i\le 8\). Block norm row sums give

% DISPLAY b.9
% SOURCE ||K−I|| ≤ p(κ),        p(t)=2t+29t²,
% SOURCE 0≤G−I,   ||G−I||≤g(κ), g(t)=8t²/(1−t)².
\begin{gather*}
\begin{aligned}
& \Vert K-I\Vert  \le  p( \kappa  ), \quad  p(t)=2t+29t^{2}, \\
& 0\le G-I, \quad  \Vert G-I\Vert \le g( \kappa  ), g(t)=8t^{2}/(1-t)^{2}.
\end{aligned}\notag
\end{gather*}

For \(K\), a diagonal block has \(2T_{9}\) plus eight squares, and each
of its seven off-diagonal blocks is \(T_{i} T_{j}-2T_{j} T_{i}\). This
explains the coefficient \(29\). Consequently

% DISPLAY b.10
% SOURCE K≥0.9771 I,     G≥I,     B≥0.9771 I.
\begin{gather*}
\begin{aligned}
& K\ge 0.9771 I, \quad  G\ge I, \quad  B\ge 0.9771 I.
\end{aligned}\notag
\end{gather*}

Also
\(\Vert D- \Gamma _{9} \otimes  I\Vert \le \sum _{i}\Vert T_{i}\Vert \le 9 \kappa\),
so \(\min \vert D\vert \ge 0.91\). Continuous scaling of \(T\) to zero
preserves its signature. \(H\ge Q_{9}^{2}\ge 0.9801\) I. Restoring \(r\)
multiplies \(K,B,H,D^{2}\) by \(r^{2}\) and \(D\) by \(r\).

\subsection{4. The real-trace two-commutator
lemma}\label{b-the-real-trace-two-commutator-lemma}

Let \(X_{1},\ldots ,X_{s}\) be Hermitian matrices with
\(\Vert X_{i}\Vert \le  \kappa\). Set
\(F^{2}=\sum _{i<j}\Vert [X_{i},X_{j}]\Vert _{F}^{2}\). For any words
\(w\) and w\textquotesingle{} of length \(q\) having the same multiset
of letters,

% DISPLAY b.11
% SOURCE |Re Tr w−Re Tr w'|
% SOURCE    ≤ ((q)_4/8) κ^(q−4) F²,                  q≥4,          (2)
\begin{gather*}
\begin{aligned}
& \vert \operatorname{Re}  \operatorname{Tr}  w-\operatorname{Re}  \operatorname{Tr}  w'\vert \\
& \le  ((q)_{4}/8) \kappa ^{q-4} F^{2}, \quad  q\ge 4,
\end{aligned}\tag{2}
\end{gather*}

where \((q)_{4}=q(q-1)(q-2)(q-3)\). The left side is zero for
\(q\le 3\).

Proof. At most \(q(q-1)/2\) adjacent swaps transform one ordering into
another. The trace difference from one swap is
\(\operatorname{Tr}  U[X_{i},X_{j}]V\). By cyclicity this is Tr \(C W\),
with \(C=[X_{i},X_{j}]\) anti-Hermitian and \(W=VU\) a word of length
\(q-2\). Exactly,

% DISPLAY b.12
% SOURCE Re Tr(CW)=(1/2)Tr(C(W−W*)).                              (3)
\begin{gather*}
\begin{aligned}
& \operatorname{Re}  \operatorname{Tr} (CW)=(1/2)\operatorname{Tr} (C(W-W^{*})).
\end{aligned}\tag{3}
\end{gather*}

Reversing \(W\) takes at most \((q-2)(q-3)/2\) adjacent swaps. Each
resulting term contains a second commutator and \(q-4\) other letters.
Hilbert--Schmidt Cauchy--Schwarz bounds its trace by
\(\kappa ^{q-4}F^{2}\). Equation (3), followed by the first sequence of
swaps, proves (2). If \(q\le 3, W\) is constant or Hermitian and (3) is
zero.

In particular, take a real-coefficient homogeneous noncommutative trace
polynomial whose commutative polynomial is zero. Sort every word to the
same canonical word for its multiset. The sorted contributions cancel
coefficient by coefficient; (2) bounds the original real trace by the
sum of absolute coefficients times \((q)_{4} \kappa ^{q-4}F^{2}/8\).
This is an explicit algebraic estimate; it does not divide by an ideal
on a singular algebraic variety.

\subsection{5. Bosonic functional calculus and an explicit coefficient
majorant}\label{b-bosonic-functional-calculus-and-an-explicit-coefficient-majorant}

Continue at \(r=1\). Put \(L=GK-I\). Since GK is similar to \(B\),

% DISPLAY b.13
% SOURCE Tr sqrt(B)=Tr sqrt(I+L)
% SOURCE           =8d+sum_(j≥1) binom(1/2,j) Tr L^j.            (4)
\begin{gather*}
\begin{aligned}
& \operatorname{Tr}  \sqrt{B}=\operatorname{Tr}  \sqrt{I+L} \\
& =8d+\sum _{j\ge 1} \binom{1/2}{j} \operatorname{Tr}  L^{j}.
\end{aligned}\tag{4}
\end{gather*}

The binomial series converges in operator norm:
\(\Vert L\Vert \le l( \kappa  )<1\), where

% DISPLAY b.14
% SOURCE l(t)=p(t)+g(t)(1+p(t)).
\begin{gather*}
\begin{aligned}
& l(t)=p(t)+g(t)(1+p(t)).
\end{aligned}\notag
\end{gather*}

The expansion \(Q_{9}^{-2}=\sum _{j\ge 0}(-1)^{j}(j+1)T_{9}^{j}\)
converges as well. Thus (4), and the corresponding expansion of 8Tr
\(\sqrt{H}\), define absolutely convergent noncommutative power series
with real coefficients.

Here is a coefficientwise majorant, not just a bound on evaluations.
Replace every letter \(T_{i}\) by a common nonnegative variable \(t\),
every scalar coefficient by its absolute value, and take maximum block
row sums before taking the eight spatial diagonal traces. The series for
\(K-I\) is majorized by p(t), for \(G-I\) by g(t), for \(L\) by l(t),
and for \(H-I\) by

% DISPLAY b.15
% SOURCE h(t)=2t+9t².
\begin{gather*}
\begin{aligned}
& h(t)=2t+9t^{2}.
\end{aligned}\notag
\end{gather*}

Since \(\sum _{j\ge 1}\vert \binom{1/2}{j}\vert z^{j}=1-\sqrt{1-z}\),
the sum of absolute trace-word coefficients of

% DISPLAY b.16
% SOURCE F_B(T)=Tr sqrt(B)−8Tr sqrt(H)
\begin{gather*}
\begin{aligned}
& F_{B}(T)=\operatorname{Tr}  \sqrt{B}-8\operatorname{Tr}  \sqrt{H}
\end{aligned}\notag
\end{gather*}

is majorized degree by degree by

% DISPLAY b.17
% SOURCE A(t)=8[2−sqrt(1−l(t))−sqrt(1−h(t))].                     (5)
\begin{gather*}
\begin{aligned}
& A(t)=8[2-\sqrt{1-l(t)}-\sqrt{1-h(t)}].
\end{aligned}\tag{5}
\end{gather*}

The eight is the spatial trace multiplicity; there is no hidden factor
\(d\), because the color trace is retained inside each word and
estimated by (2).

For scalar commuting variables, \(G K=H I_{8}\) exactly, by the rank-one
identity

% DISPLAY b.18
% SOURCE (I+vv*/q_9²)(H I−vv*)=H I,      H=q_9²+v*v.
\begin{gather*}
\begin{aligned}
& (I+vv^{*}/q_{9}^{2})(H I-vv^{*})=H I, \quad  H=q_{9}^{2}+v^{*}v.
\end{aligned}\notag
\end{gather*}

It follows that the commutative specialization of \(F_{B}\) vanishes
identically as a convergent power series. Each homogeneous polynomial
therefore has zero commutative polynomial. Apply (2) termwise, which is
allowed by the absolute majorant.

For a wholly explicit bound use \(t_{0}=1/10\). Direct arithmetic gives

% DISPLAY b.19
% SOURCE p(t_0)=49/100,
% SOURCE g(t_0)=8/81,
% SOURCE l(t_0)=5161/8100<1,
% SOURCE h(t_0)=29/100,
% SOURCE A(t_0)<5.
\begin{gather*}
\begin{aligned}
& p(t_{0})=49/100, \\
& g(t_{0})=8/81, \\
& l(t_{0})=5161/8100<1, \\
& h(t_{0})=29/100, \\
& A(t_{0})<5.
\end{aligned}\notag
\end{gather*}

If \(A(t)=\sum  a_{q} t^{q}\), all \(a_{q}\) are nonnegative and

% DISPLAY b.20
% SOURCE sum_(q≥4) a_q (q)_4 κ^(q−4)
% SOURCE    ≤24 t_0^(-4) A(t_0),             κ≤1/100.
\begin{gather*}
\begin{aligned}
& \sum _{q\ge 4} a_{q} (q)_{4} \kappa ^{q-4} \\
& \le 24 t_{0}^{-4} A(t_{0}), \quad  \kappa  \le 1/100.
\end{aligned}\notag
\end{gather*}

Indeed \((q)_{4}( \kappa  /t_{0})^{q-4}\) is maximized at \(q=4\): the
consecutive-term ratio is at most \(1/2\). Combining with (2) gives

% DISPLAY b.21
% SOURCE |Tr sqrt(B)−8Tr sqrt(H)|
% SOURCE     ≤150000 sum_(i<j)||C_ij||F².                        (6)
\begin{gather*}
\begin{aligned}
& \vert \operatorname{Tr}  \sqrt{B}-8\operatorname{Tr}  \sqrt{H}\vert \\
& \le 150000 \sum _{i<j}\Vert C_{ij}\Vert _{F}^{2}.
\end{aligned}\tag{6}
\end{gather*}

Restoring \(r\) adds \(r^{-3}\). This establishes the needed bosonic
quadratic-commutator estimate without a claim that vanishing on
commuting matrices alone implies it.

\subsection{6. Fermionic resolvent
cancellation}\label{b-fermionic-resolvent-cancellation}

Exactly,

% DISPLAY b.22
% SOURCE D²=I_16⊗H+S,
% SOURCE S=sum_(i<j)Γ_iΓ_j⊗C_ij.
\begin{gather*}
\begin{aligned}
& D^{2}=I_{16} \otimes  H+S, \\
& S=\sum _{i<j} \Gamma _{i} \Gamma _{j} \otimes  C_{ij}.
\end{aligned}\notag
\end{gather*}

The Clifford trace gives

% DISPLAY b.23
% SOURCE Tr[(I⊗H)^(-1/2)S]=0,
% SOURCE ||S||F²=16 sum_(i<j)||C_ij||F².                          (7)
\begin{gather*}
\begin{aligned}
& \operatorname{Tr} [(I \otimes  H)^{-1/2}S]=0, \\
& \Vert S\Vert _{F}^{2}=16 \sum _{i<j}\Vert C_{ij}\Vert _{F}^{2}.
\end{aligned}\tag{7}
\end{gather*}

For positive matrices \(A,A+S\ge  \lambda  I\), the square-root
resolvent identity gives

% DISPLAY b.24
% SOURCE Tr sqrt(A+S)−Tr sqrt(A)−(1/2)Tr(A^(-1/2)S)
% SOURCE   =−(1/π)integral_0^infinity t^(1/2)
% SOURCE      Tr[(A+t)^(-1)S(A+t)^(-1)S(A+S+t)^(-1)] dt.
\begin{gather*}
\begin{aligned}
& \operatorname{Tr}  \sqrt{A+S}-\operatorname{Tr}  \sqrt{A}-(1/2)\operatorname{Tr} (A^{-1/2}S) \\
& =-(1/ \pi  )\int _{0}^{\infty} t^{1/2} \\
& \operatorname{Tr} [(A+t)^{-1}S(A+t)^{-1}S(A+S+t)^{-1}] dt.
\end{aligned}\notag
\end{gather*}

The absolute value is bounded by

% DISPLAY b.25
% SOURCE ||S||F²/(8λ^(3/2)),
\begin{gather*}
\begin{aligned}
& \Vert S\Vert _{F}^{2}/(8 \lambda ^{3/2}),
\end{aligned}\notag
\end{gather*}

using
\(\int _{0}^{\infty} t^{1/2}( \lambda  +t)^{-3}dt= \pi  /(8 \lambda ^{3/2})\).
Take \(A=I \otimes  H\) and \(\lambda  =(0.91r)^{2}\). Equations (7)
prove

% DISPLAY b.26
% SOURCE |(1/2)Tr |D|−8Tr sqrt(H)|
% SOURCE    ≤(0.91)^(-3) r^(-3)sum_(i<j)||C_ij||F²
% SOURCE    <2 r^(-3)sum_(i<j)||C_ij||F².                        (8)
\begin{gather*}
\begin{aligned}
& \vert (1/2)\operatorname{Tr}  \vert D\vert -8\operatorname{Tr}  \sqrt{H}\vert \\
& \le (0.91)^{-3} r^{-3}\sum _{i<j}\Vert C_{ij}\Vert _{F}^{2} \\
& <2 r^{-3}\sum _{i<j}\Vert C_{ij}\Vert _{F}^{2}.
\end{aligned}\tag{8}
\end{gather*}

Combining (6) and (8) proves (1). The first-order fermionic term cancels
exactly under the Clifford trace; the bosonic first-order-in-commutators
term cancels by (3) and commutative specialization.

\subsection{7. Endpoint commutators, pair weights, and
absorption}\label{b-endpoint-commutators-pair-weights-and-absorption}

In a matrix representation of \(\operatorname{Hom} (C^{m},C^{n})\),

% DISPLAY b.27
% SOURCE T_i=A_i⊗I_m−I_n⊗B_i^T,
% SOURCE [T_i,T_j]=[A_i,A_j]⊗I_m−I_n⊗[B_i,B_j]^T.
\begin{gather*}
\begin{aligned}
& T_{i}=A_{i} \otimes  I_{m}-I_{n} \otimes  B_{i}^{T}, \\
& [T_{i},T_{j}]=[A_{i},A_{j}] \otimes  I_{m}-I_{n} \otimes  [B_{i},B_{j}]^{T}.
\end{aligned}\notag
\end{gather*}

There is no mixed term \([A_{i},B_{j}]\): the two endpoint actions
commute. Since traces of commutators vanish, exactly

% DISPLAY b.28
% SOURCE ||[T_i,T_j]||F²
% SOURCE    =m||[A_i,A_j]||F²+n||[B_i,B_j]||F².                 (9)
\begin{gather*}
\begin{aligned}
& \Vert [T_{i},T_{j}]\Vert _{F}^{2} \\
& =m\Vert [A_{i},A_{j}]\Vert _{F}^{2}+n\Vert [B_{i},B_{j}]\Vert _{F}^{2}.
\end{aligned}\tag{9}
\end{gather*}

For a partition with block sizes \(n_{a}\) and separations \(r_{ab}\),
summing (1) gives the pairwise-relative statement

% DISPLAY b.29
% SOURCE |sum_(a<b) E_0,ab|
% SOURCE   ≤C (R/ell³)sum_a F_a² sum_(b≠a)n_b r_ab^(-3),
% SOURCE F_a²=sum_(i<j)||[A_i^(a),A_j^(a)]||F²,   C=150002.      (10)
\begin{gather*}
\begin{aligned}
& \vert \sum _{a<b} E_{0,ab}\vert \\
& \le C (R/\ell ^{3})\sum _{a} F_{a}^{2} \sum _{b\ne a}n_{b} r_{ab}^{-3}, \\
& F_{a}^{2}=\sum _{i<j}\Vert [A_{i}^{(a)},A_{j}^{(a)}]\Vert _{F}^{2}, \quad  C=150002.
\end{aligned}\tag{10}
\end{gather*}

The sum of absolute pair defects obeys the same bound. Defining the
internal bosonic potential \(V_{a}=R F_{a}^{2}/(2\ell ^{6})\), its
multiplier is

% DISPLAY b.30
% SOURCE 2C ell³ sum_(b≠a)n_b r_ab^(-3).                         (11)
\begin{gather*}
\begin{aligned}
& 2C \ell ^{3} \sum _{b\ne a}n_{b} r_{ab}^{-3}.
\end{aligned}\tag{11}
\end{gather*}

This tends to zero once all retained pair separations tend to
\(\infty\). Its constants have no dependence on a ratio involving the
smallest unrelated center gap. At fixed finite \(N, r_{ab}\ge r_{*}\)
gives the explicit upper bound 2C \(\ell ^{3}(N-n_{a})r_{*}^{-3}\).

For clarity about full-form absorption, \(V_{a}\) by itself is not
bounded by \(H_{a}\) with coefficient one without a Yukawa remainder. If
the elementary finite-Clifford estimate is written as

% DISPLAY b.31
% SOURCE ||Y_a(A)||≤c_(n_a) (R/ell³) alpha_a,
% SOURCE alpha_a²=sum_i ||A_i^(a)||F²,
\begin{gather*}
\begin{aligned}
& \Vert Y_{a}(A)\Vert \le c_{n_{a}} (R/\ell ^{3}) \alpha _{a}, \\
& \alpha _{a}^{2}=\sum _{i} \Vert A_{i}^{(a)}\Vert _{F}^{2},
\end{aligned}\notag
\end{gather*}

then \(V_{a}\le H_{a}+c_{n_{a}}(R/\ell ^{3})\alpha _{a}\), since the
internal kinetic form is nonnegative. One explicit choice in these
conventions is \(c_{n}=24(n^{2}-1)\), with \(c_{1}=0\). Indeed, on the
\(16(n^{2}-1)\) real Majoranas, the internal mass is
\(D_{\operatorname{ad}}=\sum _{i} \Gamma _{i} \otimes  \operatorname{ad} (A_{i})\),
and \(Y_{a}=(R/(2\ell ^{3}))\theta ^{T} D_{\operatorname{ad}} \theta\).
In a real Hermitian generator basis \(D_{\operatorname{ad}}\) is purely
imaginary antisymmetric. Pairing its \(\pm  \lambda\) eigenvalues gives
\(\Vert Y_{a}\Vert =(R/(4\ell ^{3}))\operatorname{Tr} \vert D_{\operatorname{ad}}\vert\),
while

% DISPLAY b.32
% SOURCE ||D_ad||≤2sum_i||A_i||op≤6alpha_a.
\begin{gather*}
\begin{aligned}
& \Vert D_{\operatorname{ad}}\Vert \le 2\sum _{i}\Vert A_{i}\Vert _{\mathrm{op}}\le 6\alpha _{a}.
\end{aligned}\notag
\end{gather*}

This proves the stated \(c_{n}\) without an internal spectral gap. Hence
(10) is bounded by an arbitrarily small multiple of \(\sum _{a} H_{a}\),
after taking \(r_{*}\) large, plus

% DISPLAY b.33
% SOURCE 2C R sum_a c_(n_a)alpha_a sum_(b≠a)n_b r_ab^(-3).
\begin{gather*}
\begin{aligned}
& 2C R \sum _{a} c_{n_{a}}\alpha _{a} \sum _{b\ne a}n_{b} r_{ab}^{-3}.
\end{aligned}\notag
\end{gather*}

Under the endpoint condition \(\alpha _{a}\le  \kappa  r_{ab}\), this
last expression is at most

% DISPLAY b.34
% SOURCE 2C R κ sum_(a<b)(c_(n_a)n_b+c_(n_b)n_a)r_ab^(-2).
\begin{gather*}
\begin{aligned}
& 2C R \kappa  \sum _{a<b}(c_{n_{a}}n_{b}+c_{n_{b}}n_{a})r_{ab}^{-2}.
\end{aligned}\notag
\end{gather*}

It is an arbitrarily small center-Hardy cost by reducing \(\kappa\)
before choosing the final exterior radius. This is a precise way to use
the defect; replacing \(V_{a}\) by \(H_{a}\) with no remainder would be
incorrect.

For an explicit finite-N allocation, write
\(T_{\mathrm{cent}}=(R/2)t_{\mathrm{cent}}\) in the mass metric
\(\sum _{a} n_{a}\vert dz_{a}\vert ^{2}\). The nine-dimensional relative
Hardy inequality gives

% DISPLAY b.35
% SOURCE R integral r_ab^(-2)|f|²
% SOURCE   ≤(8/49)(1/n_a+1/n_b)^(-1) T_cent[f].
\begin{gather*}
\begin{aligned}
& R \int  r_{ab}^{-2}\vert f\vert ^{2} \\
& \le (8/49)(1/n_{a}+1/n_{b})^{-1} T_{\mathrm{cent}}[f].
\end{aligned}\notag
\end{gather*}

The preceding residual multiplier is therefore bounded by

% DISPLAY b.36
% SOURCE (16Cκ/49) sum_(a<b)
% SOURCE   (c_(n_a)n_b+c_(n_b)n_a)(1/n_a+1/n_b)^(-1) T_cent[f]
% SOURCE   ≤(192C/49)N^4 κ T_cent[f].
\begin{gather*}
\begin{aligned}
& (16C \kappa  /49) \sum _{a<b} \\
& (c_{n_{a}}n_{b}+c_{n_{b}}n_{a})(1/n_{a}+1/n_{b})^{-1} T_{\mathrm{cent}}[f] \\
& \le (192C/49)N^{4} \kappa  T_{\mathrm{cent}}[f].
\end{aligned}\notag
\end{gather*}

The last bound uses \(c_{n}=24(n^{2}-1)\) and
\(\sum _{a<b}n_{a}^{2}n_{b}^{2}\le N^{4}/2\). Thus, for any
\(\epsilon  >0\), the choices

% DISPLAY b.37
% SOURCE κ≤min(1/200, 49ε/(192C N^4)),
% SOURCE r_*≥[2C ell³ N/ε]^(1/3),
% SOURCE alpha_a≤κ r_ab for every incident pair,
% SOURCE r_ab≥r_* for every pair,
\begin{gather*}
\begin{aligned}
& \kappa  \le \min (1/200, 49 \epsilon  /(192C N^{4})), \\
& r_{*}\ge [2C \ell ^{3} N/ \epsilon  ]^{1/3}, \\
& \alpha _{a}\le  \kappa  r_{ab} \text{ for every incident pair }, \\
& r_{ab}\ge r_{*} \text{ for every pair },
\end{aligned}\notag
\end{gather*}

give the form bound

% DISPLAY b.38
% SOURCE integral sum_(a<b)|E_0,ab| |f|²
% SOURCE   ≤ε sum_a H_a[f]+ε T_cent[f].                         (12)
\begin{gather*}
\begin{aligned}
& \int  \sum _{a<b}\vert E_{0,ab}\vert  \vert f\vert ^{2} \\
& \le  \epsilon  \sum _{a} H_{a}[f]+ \epsilon  T_{\mathrm{cent}}[f].
\end{aligned}\tag{12}
\end{gather*}

Here the standard nonnegative internal supersymmetric forms are used on
their common core, and the center Hardy estimate is valid also for
Hilbert-valued coefficients and unitary connections. The multiplier of
each \(H_{a}\) depends only on the centers, so no derivative of an
internal-dependent weight has been omitted. Equation (12) is a bound on
this defect multiplier, not a claim that the full Hamiltonian has
already been compared with the slow forms on its right.
